\section{A proper volume obstruction}\label{sec:proper-volume}

We use one consequence of the Kashiwara--Malgrange filtration. It applies
to critical curves as well as isolated critical points.

\begin{lemma}\label{lem:proper-volume}
Let \(f:Y\to\mathbb A^1\) be a proper nonconstant morphism from a smooth
irreducible complex algebraic variety, and let \(\omega\) be a global
regular top form on \(Y\). If \(\gamma\) is a fixed compact cycle of real
dimension \(\dim_{\mathbb C}Y\), then
\[
 \int_\gamma e^{-\tau f}\omega
\]
is not a nonzero constant.
\end{lemma}

\begin{proof}
Write \(t\) for the coordinate on the target, factor \(f\) as its graph
\(i:Y\hookrightarrow Y\times\mathbb A^1_t\) followed by the projection
\(\pi\), and put \(\mathcal B_f=i_+\mathcal O_Y\).
We use the decreasing filtration \(V^\alpha\), normalized by the
nilpotence of \(\partial_t t-\alpha\) on \(\operatorname{gr}_V^\alpha\).
Negativity of the local Bernstein--Sato roots
\cite[Corollary~5.2]{Kas76} gives
\begin{equation}\label{eq:geo-star}
\delta(t-f)\in V^{>0}\mathcal B_f
 \qquad\text{\cite[Corollary~5.15 and Example~5.19]{CDM24}}.
\end{equation}
Nonconstancy ensures that every local germ of \(f\) is nonzero; no
isolated-critical-point hypothesis is used.

Let \(M=\mathcal H^0f_+\mathcal O_Y\), and let \(m\) be the image of
\(\omega\). This is a global top-form class in the relative de Rham
\emph{hypercohomology} computing \(M\), not a quotient of global forms
by global exact forms. Near the compact fiber \(f^{-1}(0)\), finitely
many local bounds in \eqref{eq:geo-star} give a common \(\alpha>0\) for
\(\omega\otimes\delta(t-f)\). Properness allows the target neighborhood
to be shrunk so that this finite cover suffices. The projection \(\pi\)
is proper on \(\operatorname{Supp}\mathcal B_f\).
For clarity, \(\operatorname{DR}_{/\mathbb A^1}\) denotes the relative
complex \(\Omega_Y^\bullet\otimes\mathcal B_f\), placed in degrees
\(-\dim Y,\ldots,0\), with differential
\(d_Y-df\wedge\partial_t\). The proper-on-support assertion of
\cite[Theorem~8.1]{CDM24} says that the image of
\[
 \mathcal H^0R\pi_*\operatorname{DR}_{/\mathbb A^1}
       (V^\alpha\mathcal B_f)\longrightarrow M
\]
is \(V^\alpha M\). The top-degree cocycle
\(\omega\otimes\delta(t-f)\) lies in this filtered complex, so
\begin{equation}\label{eq:geo-starstar}
m\in V^\alpha M\quad\text{for some }\alpha>0.
\end{equation}
This tracks the given form itself, with no parameter-dependent change
of seed.

Take global sections on the affine target and Fourier transform by
\(t\mapsto-\partial_\tau\), \(\partial_t\mapsto\tau\).
The transformed relative complex is the twisted de Rham complex
\[
 (\Omega_Y^\bullet[\tau],\ d_Y-\tau\,df\wedge)[\dim Y],
 \qquad \partial_\tau\text{ acts as }\partial_\tau-f.
\]
Analytic de Rham integration pairs its hypercohomology with compact
cycles, also for nonaffine \(Y\). Indeed,
\(d(e^{-\tau f}\beta)=e^{-\tau f}(d\beta-\tau df\wedge\beta)\),
so Stokes kills exact classes; differentiation respects the displayed
\(\partial_\tau\)-action. The period pairing therefore respects every
Weyl-algebra relation of the Fourier image \(\widehat m\).
If its period were a nonzero constant, after scaling the pairing the
cyclic submodule
\[
 \widehat N=\mathbb C[\tau]\langle\partial_\tau\rangle\widehat m
 \subset\widehat M
\]
would have a quotient to \(\mathcal O_{\mathbb A^1_\tau}\) sending
\(\widehat m\) to \(1\).
Inverse Fourier transform gives a submodule \(N\subset M\) and a
quotient \(N\to\delta_0=D_{\mathbb A^1_t}/D_{\mathbb A^1_t}t\)
sending \(m\) to its nonzero generator.
But submodules carry the induced \(V\)-filtration and morphisms
preserve it, whereas \(V^\alpha\delta_0=0\) for \(\alpha>0\)
\cite[Proposition~3.19(ii), Corollary~3.20, Example~3.18]{CDM24}.
This contradicts \eqref{eq:geo-starstar}. The cyclic quotient was constructed before
localization, so no assertion about a localized seed is needed.
\end{proof}

\section{Reduction to a constant profile}\label{sec:constant-profile}

Put \(X=\{t^2+2xy=1\}\) and \(A=\mathbb C[x,y,t]/(t^2+2xy-1)\).
The normalized residue volume \(\Omega\) on \(X\) has
\(\int_\Sigma\Omega=1\) on the real unit sphere in Cartesian
coordinates. Thus \(L(a)=\int_\Sigma a\Omega\).
In a circular chart it is a nonzero constant multiple of
\(dy\wedge dt/(2y)\). Constant factors will be suppressed in the
local regularity calculations below.

\begin{proposition}\label{prop:constant-profile}
If \(p\in A_{\leq3}\) satisfies \(L(p^m)=0\) for every \(m\geq1\),
then in some orthogonal coordinate frame
\[
 p=ay+G(x,t),\qquad \deg G\leq3,\qquad a\,[t^3]G=0.
\]
\end{proposition}

\begin{proof}
Choose a representative \(P=P_d+\cdots+P_0\) of minimal degree \(d\),
with homogeneous pieces \(P_j\). The cases \(d\leq1\) already have
the stated form. For \(d\geq2\), minimality implies
\(P_d|_{\{q=0\}}\neq0\): otherwise \(P_d=qH\), and subtracting
\((q-1)H\) lowers the degree.

Compactify \(X\) by
\[
 \bar X=\{T^2+2XY=R^2\}\subset\mathbb P^3.
\]
The actual closure \(\Gamma_d\) of the graph of \(p\) is proper over
the phase line and is given by
\begin{equation}\label{eq:geo-1}
\sum_{j=0}^d R^{d-j}P_j(X,Y,T)-sR^d=0.
\end{equation}
There is no component supported on \(R=0\), because the nonzero
restriction of \(P_d\) has only finitely many zeros on the boundary
conic. All modifications below are strict transforms of this
closure, and are isomorphisms over \(X\).

Move a boundary point to \([X:Y:T:R]=[0:1:0:0]\) by an orthogonal
change of frame, and use coordinates
\begin{equation}\label{eq:geo-2}
(X,Y,T,R)=((z^2-u^2)/2,1,u,z).
\end{equation}
The original volume is, up to a nonzero constant,
\begin{equation}\label{eq:geo-3}
\Omega=\frac{du\wedge dz}{2z^2}.
\end{equation}

Suppose first that \(d=2\).
At a boundary root, the coefficient of \(y^2\) in \(P_2\) is zero.
If the root of \(P_2|_{\{q=0\}}\) is repeated, its derivative in
\(u\) also vanishes, so the coefficient of \(yt\) is zero.
Replacing \(xy\) by \((1-t^2)/2\) then gives
\(p=ay+G(x,t)\), as required.
If no such frame exists, all boundary roots are simple.
The local equation \(F=0\) in \eqref{eq:geo-1} is then smooth at every boundary
point, with \(F_s=-z^2\). Its regular dualizing residue \(\eta_F\)
satisfies
\[
 \Omega=-\tfrac12\eta_F.
\]
Consequently \(\Omega\) extends to the whole smooth proper graph.
Lemma~\ref{lem:proper-volume} contradicts
\(\int_\Sigma e^{-\tau p}\Omega=1\).
This settles \(d=2\), without a separate quadratic classification.

Now let \(d=3\), and suppose that no frame satisfies the conclusion.
At a boundary point in \eqref{eq:geo-2}, its equation has derivatives
\[
 F_u(0,0,s)=[y^2t]P_3,\qquad
 F_z(0,0,s)=[y^2]P_2,
\]
while the root condition is \([y^3]P_3=0\).
If either derivative is nonzero, the graph is smooth there.
Since \(F_s=-z^3\), its dualizing residue satisfies
\begin{equation}\label{eq:geo-4}
\Omega=-\tfrac z2\,\eta_F,
\end{equation}
which extends regularly across that point.

It remains to treat points where both derivatives vanish.
If \(C=[xy^2]P_3\), replace \(P\) by
\(P-\frac C2y(q-1)\). This removes \(xy^2\) and changes neither
\(p\) nor the graph on \(\bar X\). The representative is now affine
linear in \(y\). Eliminating its remaining \(xy\) factors gives
\[
 p=y(a+bt+ct^2)+G(x,t),\qquad \deg G\leq3.
\]
In \eqref{eq:geo-2} the graph equation is
\begin{equation}\label{eq:geo-5}
F=cu^2+buz+az^2+R_{\geq3}(u,z)-sz^3=0.
\end{equation}
Indeed \(x^it^j\) in \(G\) contributes
\(2^{-i}(z^2-u^2)^iu^jz^{3-i-j}\), of total order at least \(3+i\).
A constant profile with \(a[t^3]G=0\) would finish the proof. We resolve the other cases.

If \(b^2-4ac\neq0\), blow up the boundary section.
In the chart \(u=zv\), the strict transform is \(F=z^2F_1\), with
\[
 F_1=cv^2+bv+a+zK(v,z,s),\qquad (F_1)_s=-z.
\]
Its exceptional intersections are the simple roots of the profile,
so it is smooth there, and
\begin{equation}\label{eq:geo-6}
\Omega=\frac{dv\wedge dz}{2z}=-\tfrac12\eta_{F_1}.
\end{equation}
If \(c=0,b\neq0\), the additional tangent direction is covered by
\(z=uv\). Here \(F=u^2F_2\), and the exceptional point \(u=v=0\)
has \((F_2)_v=b\neq0\). Moreover
\begin{equation}\label{eq:geo-7}
(F_2)_s=-uv^3,\qquad
 \Omega=\frac{du\wedge dv}{2uv^2}=-\tfrac v2\,\eta_{F_2}.
\end{equation}
Thus both charts are smooth along their full exceptional boundary
and the volume is regular there.

The remaining nonconstant profile has \(c\neq0\) and
\(b^2-4ac=0\). Divide the local equation by \(c\) and rescale the phase coordinate, so that the profile is \((t-r)^2\),
and put \(w=u-rz\). Blow up in the chart \(w=zv\), which contains
the entire exceptional intersection. The strict equation is
\begin{equation}\label{eq:geo-8}
F_1=v^2+z(k(v)-s)+z^2\rho(v,z)=0.
\end{equation}
It is smooth on the boundary except possibly at
\((v,z,s)=(0,0,k(0))\).
The Hessian determinant at that point is \(-2\), independently of
\(k\) and \(\rho\); it is an ordinary \(A_1\) surface double point.
By \eqref{eq:geo-6}, \(\Omega\) is a constant multiple of its dualizing generator.
Blow up this point. The strict transform is smooth, with a smooth
conic exceptional divisor, and the resolution is crepant:
the ambient point blow-up adds twice the exceptional divisor to
the canonical divisor, while the double-point hypersurface has
multiplicity two. Adjunction cancels these contributions.
The given volume therefore extends holomorphically over the
resolution.


For a constant profile, our contrary assumption says \(a\neq0\) and
\(h=[t^3]G\neq0\). Write
\(G(0,t)=ht^3+g_2t^2+g_1t+g_0\).
Equation \eqref{eq:geo-5} is then
\begin{equation}\label{eq:geo-9}
F=az^2+hu^3+g_2u^2z+g_1uz^2+(g_0-s)z^3+R_{\geq4}(u,z).
\end{equation}
Blow up in the chart \(z=uv\). The strict transform is
\begin{equation}\label{eq:geo-10}
F_1=av^2+u\bigl(h+g_2v+g_1v^2+(g_0-s)v^3\bigr)
          +u^2H(u,v)=0.
\end{equation}
The exceptional intersection is the single point \(u=v=0\), where
\((F_1)_u=h\neq0\). Hence it is smooth for every finite \(s\).
Furthermore \((F_1)_s=-uv^3\), and \eqref{eq:geo-7} again gives
\(\Omega=-v\eta_{F_1}/2\), a regular form.
Equivalently, solving \(u=u(v,s)\) gives
\[
 \Omega=\pm\frac{v}{2(F_1)_u}\,dv\wedge ds.
\]
The complementary chart \(u=zv\) has strict equation \(a+O(z)=0\)
and therefore no exceptional intersection. These charts cover the
complete exceptional divisor; no boundary-supported component
has been retained.

There are finitely many boundary points. The indicated blow-ups
of boundary sections, followed by the finitely many ordinary
double-point resolutions, give a smooth irreducible surface \(Y\)
proper over \(\mathbb A^1_s\), containing \(X\), on which the original
volume \(\Omega\) is regular. The phase is nonconstant and the
original sphere remains a compact cycle of nonzero volume.
Purity gives \(\int_\Sigma e^{-\tau p}\Omega=1\), contradicting
Lemma~\ref{lem:proper-volume}.
A frame with constant profile and \(a[t^3]G=0\) must therefore exist.
\end{proof}
